\begin{table}[!htbp]
\centering
\footnotesize
\setlength{\tabcolsep}{3.5pt}
\caption*{\textbf{Table C6.} Out-of-zone and in-zone early voters by counting jurisdiction (conservative valid-vote share, percentage points)}
\begin{tabularx}{\linewidth}{lCCCCCC}
\toprule
Election & Jurisdictions & Out-of-zone minus in-zone: median & Range & Largest absolute gap & In-zone stations: median absolute deviation from their jurisdiction & Largest station deviation \\
\midrule
2017 & 25 & $-$3.5 & $-$5.4 to $-$0.6 & 5.4 & 1.7 & 14.8 \\
2022 & 25 & 1.7 & $-$3.2 to 3.4 & 3.4 & 2.2 & 19.9 \\
2024 & 48 & $-$2.5 & $-$6.0 to 0.1 & 6.0 & 1.8 & 19.2 \\
2025 & 25 & 0.4 & $-$2.3 to 2.7 & 2.7 & 1.4 & 21.9 \\
2026 & 25 & 3.6 & 1.7 to 7.7 & 7.7 & 1.9 & 25.7 \\
\bottomrule
\end{tabularx}
\par\vspace{2pt}{\raggedright\footnotesize\textit{Notes.} Each jurisdiction's out-of-zone early ballots are published as one total, so the first three columns describe whole districts (constituencies in 2024). The tolerance $\delta $ bounds the same difference for each neighborhood's residents, which is not observed and may vary more within a district than these district-level gaps; the station columns show how far neighborhoods' in-zone shares depart from their district's in-zone share. These are references for $\delta $, not tests of it.\par}
\end{table}
